% r03_fill_K1: proof of the torus comparison (K_0), part I: the statement at the strength proved, the two block-spin flows.
% Input from sec/r03_rot_torus.tex after the proof of \Cref{rot:hyp:K}. Parts II-VI follow as sec/r03_fill_K2a.tex ... K6.tex.
\providecommand{\Xk}{\ensuremath{X_k}}
\providecommand{\Yk}{\ensuremath{Y_k}}
\providecommand{\hs}{\ensuremath{\hat s}}
\providecommand{\vr}{\ensuremath{\varrho}}
\providecommand{\Bdel}{\ensuremath{(\mathrm B_\Delta)}}
\providecommand{\ipp}[1]{\ensuremath{\langle #1\rangle}}

\subsection{Proof of the torus comparison}\label{ssec:rot:Kproof}

This subsection proves the one-sided bounds \hypFz\ and \hypV\ of \Cref{rot:hyp:F}, in the total form in which the vacuum constants are absent and the source-linear constants tend to zero, for the walled configurations of the class of \Cref{rot:rem:class}; together with the existence of the axis-lattice limits (\Cref{rot:thm:H6W}) and \Cref{rot:cor:walled} this completes the proof of \Cref{rot:hyp:K}. Throughout, $R\in\Rgen$ is fixed, $5R$ and $5R^{-1}$ are integer matrices (\Cref{rot:lem:geom}(a)), and we write
\[
\Xk\defeq\Lk=\eps_k\Z^3\quad\text{(the axis lattice)},\qquad\Yk\defeq R^{-1}\Lk\quad\text{(the rotated lattice)},\qquad
Y_{k-j}\defeq R^{-1}(5^j\eps_k\Z^3),
\]
so that $Y_{k-j}\subset\Xk\cap\Yk$ for $j\ge1$; the frame of $\Xk$ is $(e_\alpha)$ and that of $\Yk$ is $(R^{-1}e_\alpha)$. The tori are $T_1=\DJ=\R^3/P_J$, $P_J=50MJ\,\Z^3$, and $T_2=R^{-1}\DJ=\R^3/R^{-1}P_J$ (\Cref{rot:hyp:F}); the involution is $S=R^{-1}\sigma_\nu$ of \Cref{rot:lem:inv}. In this subsection $r$ denotes the number of block-spin steps, never the modulus $r_x=|\phi(x)|$, which does not occur.

\begin{theorem}[the torus comparison at one configuration]\label{rot:thm:CT}
Let the couplings be in King's range and $\fc=(J,\psi,n)$ a walled configuration with $J\in5^{n_*}\N$, $J\ge J_0(\delta_r)$, $n>0$, the wall condition $n\sS_{\Xk}(\psi)>\Delta_*|T_1|/2$ ($\Delta_*$ the constant of \Cref{rot:prop:CT}(b)) and block alignment (the source box and the torus unions of blocks of the final scale); here $n_*$ is the number of large-field scales of the flow of $\Yk$, fixed by the couplings, and $J_0(\delta_r)$ the non-wrapping threshold of \Cref{rot:lem:TC} at the coarsest spacing $\delta_r$ of the flow (a condition void for large $k$, since $\delta_r\to0$). Then for every real smooth $P_J$-periodic $2$-form $G$:
\begin{enumerate}[label=\textup{(\roman*)},leftmargin=*]
\item $\langle e^{zF(G)}\rangle_{\fc,\Yk}-\langle e^{zF(G)}\rangle_{\fc,\Xk}\to0$ as $k\to\infty$, locally uniformly in $z\in\C$; in particular all moments $\langle F(G)^m\rangle$ differ by $o(1)$;
\item both families converge, locally uniformly on $\C$, to the same entire function $\Wv_\fc(zG)$, and all moments converge to the same limits;
\item the same holds with $R$ replaced by any $R'\in\Rgen$.
\end{enumerate}
At $e=0$ the statement holds for every walled configuration with no condition: $\langle e^{zF(G)}\rangle_{\fc,\Lambda}=e^{z^2q_\Lambda(G)/2}$ exactly (\Cref{rot:lem:CM}), and $q_{\Xk}(G),q_{\Yk}(G)\to q(G)$ (\Cref{rot:lem:q}).
\end{theorem}

The three conditions on $\fc$ are those of \Cref{rot:rem:class}, and the configuration produced by \Cref{rot:thm:REStoTIGHT} satisfies them (its wall height may be raised at will, \Cref{rot:rem:Eproof}(iii), and $J$, the source box and the torus are chosen in \Cref{rot:thm:shell}). The proof is assembled in \Cref{rot:prop:Kforward} and the paragraph after it, from the comparison of the effective actions of two block-spin flows at equal spacing (\Cref{rot:thm:Asrc}), King's upper bound for the rotated lattice (\Cref{rot:prop:H3}), the hybrid lower bound (\Cref{rot:prop:H4}) and a polynomial lower bound at the wall (\Cref{rot:thm:FS}), closed by \Cref{rot:thm:sandwich}.

\begin{lemma}[admissibility]\label{rot:lem:adm}
\textup{(i)} $P_J\subset M\Z^3\cap R^{-1}(M\Z^3)$. \textup{(ii)} $SP_J=R^{-1}P_J\subset M\Z^3\cap R^{-1}(M\Z^3)$. \textup{(iii)} For $0\le j\le k$, $P_J\subset5^j\eps_kM\Z^3\cap R^{-1}(5^j\eps_kM\Z^3)$, and the same for $R^{-1}P_J$: every block lattice of both flows below descends to both tori.
\end{lemma}
\begin{proof}
(i) $50J\,R^{\mathrm T}=10J\cdot5R^{\mathrm T}$ is integral, so $RP_J=R^{-\mathrm T}P_J\subset10MJ\,\Z^3$, and $R^{-1}(M\Z^3)\supset P_J$. (ii) $\sigma_\nu\Z^3=\Z^3$ gives $SP_J=R^{-1}P_J$, and $R^{-1}P_J\subset10MJ\,\Z^3$ as in (i). (iii) $5^{k-j}\in\Z$.
\end{proof}

\subsubsection*{The two block-spin flows}

\begin{definition}[the standard flow of $\Yk$ and the hybrid flow of $\Xk$]\label{rot:def:flows}
Fix $k$ and $r\le k$; $\delta_r\defeq5^r\eps_k$ is the physical spacing of the coarse lattice $Y_{k-r}$. In the unit scale of level $r$ write $\eta\defeq5^{1-r}$ and $\eta'\defeq\eta/5=5^{-r}$. Blocks are centred at every level (\Cref{rot:conv:C}).
\begin{itemize}[leftmargin=*]
\item \emph{Standard flow on $\Yk$} (flow $Y$): Ba{\l}aban's and King's block-spin transformation \cite[(2.1)--(2.11)]{Ba1}, \cite[(2.10)--(2.17)]{K1} in the frame of $\Yk$: centred cubes of $125$ points, $Y_{k-j+1}\to Y_{k-j}$, with contours made of fine $\Yk$-bonds in the axial order $f_1,f_2,f_3$ of the frame inside each block.
\item \emph{Hybrid flow on $\Xk$} (flow $X$): the first step averages over the centred tiles $y+\eps B_0$, $y\in Y_{k-1}$, with
\[
B_0\defeq\Z^3\cap R^{-1}\big(-\tfrac52,\tfrac52\big)^3
\]
($125$ sites, a complete residue system of $\Z^3/5R^{-1}\Z^3$, \Cref{rot:lem:tile}), onto $Y_{k-1}$; this is \cite[(4.3)--(4.4)]{K2}, with the photon average read frame-covariantly, $(\hat QA)_\alpha(y)=\frac1{125}\sum_{x\in B(y)}f_\alpha\cdot A(x)$, $f_\alpha=R^{-1}e_\alpha$ (\Cref{rot:rem:frameav}). Steps $2,\dots,r$ are the standard flow's steps $Y_{k-1}\to Y_{k-r}$ with the same blocks, except that every rotated bond of $Y_{k-1}$ occurring in a contour is replaced by a fixed path of fine $\Xk$-bonds (\cite[p.~336]{K2}: ``succeeding steps of the renormalization transformation in the usual way, that is without rotating the blocks'').
\item \emph{Contours, rotation-paired.} Let $\Gamma^{\rm st}_{z,w}$ be King's coordinate staircase from $z$ to $w$ \cite[(2.12)]{K1} and $\Gamma^{(r)}_{y,x}$ his composite contour \cite[(2.11)]{K1}, which the standard flow uses in $Y$'s frame. The legs of a $Y$-frame staircase between two centres of $Y_{k-1}$-blocks are concatenations of the $\eta$-bonds $(u,u+5\eps f_\alpha)$, $u\in Y_{k-1}$. The hybrid composite contour from the level-$r$ block centre $y$ to $x'\in\Xk$ is the same staircase with every $\eta$-bond $(u,u+5\eps f_\alpha)$ replaced by the \emph{fixed} fine axis staircase $u+\eps\gamma_\alpha$ (and the reversed bond by its reverse), where, for $R=R_{12}(\thz)$ and $f_\alpha=R^{-1}e_\alpha$ with $5f_1=(4,3,0)$, $5f_2=(-3,4,0)$,
\[
\gamma_1:\ e_1,e_2,e_1,e_2,e_1,e_2,e_1\ \ (\text{to }(4,3,0)),\qquad
\gamma_2:\ e_2,-e_1,e_2,-e_1,e_2,-e_1,e_2\ \ (\text{to }(-3,4,0)),\qquad
\gamma_3:\ e_3^{\,5},
\]
followed, inside the first-step tile, by a fine axis path of at most $5$ bonds (\Cref{rot:lem:tile}(ii)), chosen once as a function of the displacement only and the same in every tile, so that the block and contour system of the flow is invariant under $Y_{k-1}$, in particular under the periods of both tori; this is the first-step contour of \cite[(4.4)]{K2}. (For the other generators conjugate by the permutation and sign matrices of \Cref{rot:lem:K}(c).) The gauge field along a composite contour is read by King's convention \cite[(2.19)]{K1} (\Cref{rot:conv:C}(ii)).
\end{itemize}
Both flows use King's coefficients $a_j=a(1-5^{-2})/(1-5^{-2j})$ \cite[p.~653]{K1}, take $r$ steps and land on the same coarse lattice $Y_{k-r}/P$ with the same unit-scaled fields $(A_r,\phi_r)$ and the same small-field indicator $\chi_r$ \cite[(3.2)]{K1}. The effective actions are written so that
\[
\Zun_{\Xk,T}(s)=\int(dA_r)(d\phi_r)\,e^{-\hs^{(r)}_T(A_r,\phi_r;s)},\qquad
\Zun_{\Yk,T}(s)=\int(dA_r)(d\phi_r)\,e^{-S^{(r)}_T(A_r,\phi_r;s)}
\]
\emph{exactly}, every $c$-number kept inside: this is \cite[(3.14), (3.90)]{K1} with $E_0$ inserted from \cite[(3.89)]{K1}.
\end{definition}

\begin{lemma}[the paired contours]\label{rot:lem:pair}
Every point of $u+\eps\gamma_\alpha$ lies within $0.6\eps$ of the segment $[u,u+5\eps f_\alpha]$, and the two bound a planar region of area at most $5\eps^2$. Each level-$i$ piece of the hybrid composite contour lies inside its level-$(i+1)$ block and the in-tile piece inside its tile, so the contour lies inside the level-$r$ block; in the unit scale of level $r$ its length is at most $2.1$ (the standard flow's at most $1.5$); the Hausdorff distance between the level-$i$ pieces of the two flows is at most $\eta'$ for $i\ge1$ and at most $2\sqrt{13}\,\eta'$ in the first-step cell; and the closed polygon formed by the two composite contours and the matching segment $[x',\tau x']$ (\Cref{rot:def:match}) bounds a piecewise planar surface of area at most $1.5\,\eta'+33\,\eta'^2$.
\end{lemma}
\begin{proof}
A finite check. The staircase $\gamma_1$ alternates $e_1$ and $e_2$ steps from $0$ to $(4,3,0)$ along the segment of slope $3/4$; its vertices $(1,0),(1,1),(2,1),(2,2),(3,2),(3,3),(4,3)$ have distances $\le0.6$ from the segment (the maximum, $3/5$, at $(1,0)$ and $(3,3)$), and the region between them is a union of triangles of total area $\le5$; $\gamma_2$ is $\gamma_1$ rotated by a right angle, and $\gamma_3$ is the segment itself. The centres of the blocks are at distance $\ge\frac52\eps$ from the block faces, against the deviation $0.6\eps$, which gives the containments; the lengths are $7$ fine bonds per $\eta$-bond for $\alpha=1,2$ and $5$ for $\alpha=3$, against $5$ for the rotated bond, plus the in-tile piece of at most $5$ bonds, which gives $2.1$ against $1.5$ per unit length of staircase in the unit scale; the Hausdorff distances and the area are the sums over the pieces of the two levels, the first-step cell contributing $\max_{B_0}|x|\le\sqrt{13}$ (\Cref{rot:lem:tile}(ii)).
\end{proof}

\begin{convention}[centred blocks; the contour convention]\label{rot:conv:C}
\textup{(i)} King's blocks are corner-rooted; both flows here use blocks centred at the sites of the next coarse lattice. Since $5$ is odd, centred cubes and centred tiles are inversion-symmetric. King's blocks are translates of these, and the printed estimates use only that blocks are translates of one cube and that contours stay inside blocks with bounded length (tiling, nesting, site in block with bounded diameter, all blocks translates of one cube); the four properties hold for centred blocks at block factor $5$ by exact enumeration, the fluctuation-integral estimates of \cite{Ba1,Ba2,Ba3,Ba4} being insensitive to the root of the block once the central symmetry of the fixed staircases is given. With centred blocks the inversion about any site of the coarsest lattice is an exact symmetry of both flows at every scale, which removes every first-moment term from both effective actions. \textup{(ii)} The gauge field along a composite contour is read in King's convention \cite[(2.19)]{K1}: if $A=\sum_{j<r}A'^{(j),\eps}+A^{(r),\eps}$ is the decomposition of the fine field into the lifts of the integrated fluctuations \cite[(2.18)]{K1}, then $A(\Gamma^{(r)}_{y,x})=\sum_{j=0}^{r-1}A'^{(j),\eps}(\Gamma^{(j+1)}_{x_{j+1},x})+A^{(r),\eps}(\Gamma^{(r)}_{y,x})$: the stage-$j$ fluctuation is integrated only along the part of the contour inside the $(j+1)$-block. This is what iterating the one-step transformations produces, and it is a definition in all three prints, \cite[(2.19)]{K1}, \cite[(3.34)]{Ba1}, \cite[(1.3)]{Ba3}. Under the other reading (one fine field along the whole composite contour) a staircase and the straight rotated bond with the same displacement carry different order-$e^2$ self-energies, by $0.0265\,e^2$ per unit length, an extensive amount; under \cite[(2.19)]{K1} the mismatch is $O(e^2\eps)$ per contour and its total vanishes.
\end{convention}

\begin{remark}[the photon average is frame-covariant]\label{rot:rem:frameav}
\cite[(4.3)]{K2} prints $(\hat QA_\mu)(y)=5^{-d}\sum_{x\in B(y)}A_\mu(x)$, which puts the axis component $A_\mu$ on the rotated bond direction. For a constant field $B$ it produces the coarse field $R^{-1}B$ while the fine contours see $B$, so King's constant-background step \cite[(3.43)--(3.46)]{K1} is inconsistent with it. With the frame-covariant average of \Cref{rot:def:flows}, $\hat Q^*\hat Q=\hat P\otimes\mathbf 1$ because $\sum_\alpha f_\alpha f_\alpha^{\mathrm T}=\mathbf 1$, the hybrid photon covariances are the scalar ones times the identity, and a constant background passes through every step unchanged. This repair is used throughout.
\end{remark}

\begin{lemma}[the tile average is a block-spin average]\label{rot:lem:tile}
Let $B(y)\defeq y+\eps B_0$ for $y\in Y_{k-1}$.
\begin{enumerate}[label=\textup{(\roman*)},leftmargin=*]
\item $B_0$ is a complete residue system of $\Z^3/5R^{-1}\Z^3$ with $|B_0|=125$; the tiles $B(y)$, $y\in Y_{k-1}$, partition $\Xk$, and for every period lattice $P$ of \Cref{rot:lem:adm} the tiles $B(y)$, $y\in Y_{k-1}/P$, partition $\Xk/P$.
\item $B_0=-B_0$, $\max_{x\in B_0}|x|^2=13$, $|x|_\infty\le3$ on $B_0$, and every $x\in B_0$ is joined to $0$ by a nearest-neighbour path inside $B_0$, monotone in each coordinate, of length $|x|_1\le5$.
\item With $\hat Q$ and $Q(A)$ of \Cref{rot:def:flows} and the pairings $\ip uv_\Lambda=\sum_xs^3u(x)\cdot v(x)$ of the two lattices, $\hat Q\hat Q^*=\mathbf 1$ on coarse bond fields and $Q(A)Q(A)^*=\mathbf 1$ on coarse scalar fields, for every bond field $A$. Hence $\hat Q^*\hat Q$ and $Q(A)^*Q(A)$ are orthogonal projections, the Gaussian block-spin kernels $\exp[-\frac a2\|A_1-\hat QA\|^2]$, $\exp[-\frac a2\|\phi_1-Q(A)\phi\|^2]$ have the normalisations of \cite[(2.10)--(2.14)]{K1}, and the first hybrid step is King's renormalisation transformation with $B(y)$ in place of the cube.
\item The composite of the first step with the $j-1$ standard steps that follow it is the equal-weight average over the composite blocks $B^{(j)}(y)$, $y\in Y_{k-j}$ (unions of $125^{j-1}$ tiles), which partition $\Xk$ and have diameter $\le C5^j\eps$; the composite Gaussian kernel has King's coefficient $a_j$, as for the standard flow.
\end{enumerate}
\end{lemma}
\begin{proof}
(i) The open cube $R^{-1}(-\frac52,\frac52)^3$ is a fundamental domain of the lattice $5R^{-1}\Z^3$, and no point of $\Z^3$ lies on its boundary: for $R=R_{12}(\thz)$ the boundary planes are $4x_1+3x_2=\pm\frac{25}2$, $-3x_1+4x_2=\pm\frac{25}2$, $x_3=\pm\frac52$, none of which has an integer solution. A fundamental domain whose boundary carries no lattice point contains exactly one point of every coset, so $B_0$ is a complete residue system and $|B_0|=|\det5R^{-1}|=125$. Hence every $x\in\Z^3$ is uniquely $x=5R^{-1}m+b$ with $m\in\Z^3$, $b\in B_0$, i.e.\ $\eps x=y+\eps b$ with $y=5\eps R^{-1}m\in Y_{k-1}$ unique: the tiles partition $\Xk$. For admissible $P$, $P\subset Y_{k-1}$ (\Cref{rot:lem:adm}(iii) with $j=1$), so $P$-translates of tiles are tiles and the partition descends. (ii) The cube is inversion-symmetric, so $B_0=-B_0$. The rest is a finite check over the $125$ points: $|x|^2\le13$ with equality at $(\pm3,0,\pm2)$, $(0,\pm3,\pm2)$; $|x|_\infty\le3$; for each $x$ a coordinate-monotone path from $0$ stays in $B_0$, the longest having $|x|_1=5$ steps. (iii) For a coarse bond field $\psi$ on $Y_{k-1}$ the adjoint of $\hat Q$ with respect to the pairings is $(\hat Q^*\psi)(x)=\sum_\alpha\psi_\alpha(y(x))f_\alpha$, $y(x)$ the centre of the tile containing $x$, because $(5\eps)^3/125=\eps^3$; then $(\hat Q\hat Q^*\psi)_\alpha(y)=\frac1{125}\sum_{x\in B(y)}f_\alpha\cdot\sum_\beta\psi_\beta(y)f_\beta=\psi_\alpha(y)$. For the scalar field, $(Q(A)\phi)(y)=\frac1{125}\sum_{x\in B(y)}U(A(\Gamma_{yx}))\phi(x)$ with $|U|=1$ has adjoint $(Q(A)^*\chi)(x)=\overline{U(A(\Gamma_{y(x)x}))}\chi(y(x))$, and $Q(A)Q(A)^*=\mathbf 1$. An operator with $QQ^*=\mathbf 1$ has $Q^*Q$ an orthogonal projection, and the Gaussian integrals $\int dA_1\exp[-\frac a2\|A_1-\hat QA\|^2]$, $\int d\phi_1\exp[-\frac a2\|\phi_1-Q(A)\phi\|^2]$ are constants independent of $A$ and $\phi$, which is all that \cite[(2.10)--(2.14)]{K1} uses of the block. (iv) The composition of equal-weight averages over successive partitions is the equal-weight average over the composite blocks; the diameter bound is (ii) at step $1$ and the cube's diameter at the later steps; the coefficient $a_j$ is King's computation \cite[p.~653]{K1}, which uses only the form of the kernels.
\end{proof}

\begin{lemma}[flux through the sliver]\label{rot:lem:flux}
For two lattice curves with the same endpoints, within Hausdorff distance $w$ and of length $\le\ell$, $|\int_\gamma A-\int_{\gamma'}A|\le\|dA\|_\infty\,\mathrm{area}(\Sigma)$ for a surface $\Sigma$ they bound together with the matching segment, and in particular $\le C\|dA\|_\infty w\ell$. In the unit scale of level $r$ the paired systems of \Cref{rot:def:flows} bound a surface of area $\le1.5\eta'+33\eta'^2$ (\Cref{rot:lem:pair}), so the phases $U(A(\Gamma))$ of the two flows at matched endpoints differ by at most $C\eta'\|dA\|_\infty$, a relative $\eta'$ against the Wilson-line phase.
\end{lemma}
\begin{proof}
Stokes' theorem on the lattice (the sum of the plaquette fluxes over $\Sigma$). The pairing is a condition, not a convenience: if the enclosed loops did not shrink with the scale, the flux $\sim e_rpA\ell^2$ would be of the same order as the Wilson-line phase $e_rA\ell$ at the coarse momenta $p\ell=O(1)$ at which the block-spin comparison is made, and \Cref{rot:thm:A} would fail; \cite[(4.4)]{K2} asks only that the contour ``lie entirely inside $B(y)$'' and leaves its shape open.
\end{proof}
